% ============================================================
%  Table III: P2 H5 Fixed-Secondary Timing Reference & Loss-Recovery Instances
%  ITCA / IEEEtran rule: period after caption.
% ============================================================
\begin{table}[t]
\centering
\caption{Fixed-Secondary Absolute Local Prediction Errors and Unilateral Loss-Recovery Instances (Runs 2715--2717, \texttt{time} Profile).}
\label{tab:p2_time}
\footnotesize
\setlength{\tabcolsep}{1pt}
\begin{tabular*}{\columnwidth}{@{}l@{\extracolsep{\fill}}lcccrr@{}}
\toprule
 & & & Probe & In-sched. & \multicolumn{2}{c}{$|e|$ (ms)} \\
\cmidrule(l){6-7}
Run & Side & Joint & coverage & recovery & Median & P95 \\
\midrule
2715 & BLE  & 6/6 & 24/24 & N/A (0/0) & 0.223 & 0.318 \\
     & IEEE &     & 24/24 & N/A (0/0) & 0.129 & 0.295 \\
\addlinespace
2716 & BLE  & 5/6 & 19/24 & 1/1       & 0.242 & 0.395 \\
     & IEEE &     & 24/24 & N/A (0/0) & 0.144 & 0.330 \\
\addlinespace
2717 & BLE  & 5/6 & 24/24 & N/A (0/0) & 0.222 & 0.317 \\
     & IEEE &     & 19/24 & 1/1       & 0.134 & 0.336 \\
\bottomrule
\end{tabular*}
\par\smallskip
\raggedright\scriptsize
Joint counts are per run. Run 2715 had no intervention; the BLE and IEEE feeders were disconnected in 2716 and 2717, respectively. Coverage counts held-out probes; recovery counts in-schedule events (N/A when none). Each side had one post-transmission stale event, excluded here. Across 24, 19 and 19 co-covered probes, median cross-side differences ranged from 0.140 to 0.153 ms and P95 from 0.202 to 0.264 ms. Each run had 0/6 tasks eligible for timed software processing under uncalibrated RF bounds. $|e|$ is the absolute local prediction error.
\end{table}
